\answer{ex:16:1}{血液は$\tau=t_{1/2}/\ln2=14.4$~minのRC回路である。最大と最小の比は$2^{T/t_{1/2}}=64$、基本波は$0.55$に減衰する。比が$2$になるのは$t_{1/2}=T=60$~minのとき。}
\answer{ex:16:2}{$\varphi^*=\arcsin\bigl((\omega-\omega_0)/K_\varphi\bigr)\approx41$~min（$\omega-\omega_0\approx0.047$~rad/day）、緩和時間は$1/(K_\varphi\cos\varphi^*)\approx3.9$~day。$+6$~hと$-6$~hのずれのあとは、それぞれ$7.7$~dayと$7.9$~day。}
\answer{ex:16:3}{$K_c(n)=\sec^n(\pi/n)$：$8$と$2.37$、誤差は$11\%$と$30\%$。$n\to\infty$で$K_c\to1$、誤差$\to50\%$。}
\answer{ex:16:4}{両極端の指令（アクセル$80$、ブレーキ$0$）で$0.76$~s後に$f=120$となる。「アクセルだけ」では$2.33$~sで、3倍長い。}
\answer{ex:16:5}{$3\arctan(\omega\tau)+\omega d=\pi$、$K_c=(1+\omega^2\tau^2)^{3/2}$：$K_c=8$、$3.54$、$2.50$、$1.76$、周期は$3.63$、$5.46$、$6.86$、$9.27\,\tau$。$0.9K_c$では振動が減衰し、$1.1K_c$では持続振動に落ち着く。}
\answer{ex:16:6}{$0.60\bar c$、$0.87\bar c$、$0.90\bar c$、極限は$0.91\bar c$。濃度が一定なら応答はゼロである。このような受け手は分量しか読まない。}
\answer{ex:16:7}{自由課題。フィードバックならリズムは$K>8$でしか現れず、周期は段の$\tau$に比例し、$K$にはほとんど依存しない（$K=8.5$で$3.64\,\tau$、$K=40$で$4.23\,\tau$）。$K$を1.5倍変えても周期は$2$--$4\%$しか動かず、誤差より小さい。決め手はループを開くこと（第1段の入力に最終段のホルモンを一定量与えると、この種のリズムは消える）か、周期の異なる位相で短時間ホルモンを追加したときの応答である。}
